% ══════════════════════════════════════════════════════════════════════════
%  Chapter 6: Linear Threshold Functions, Reed--Muller Codes, and
%             the Group-Theoretic Perspective
% ══════════════════════════════════════════════════════════════════════════

\chapter{Linear Threshold Functions, Codes, and Groups}\label{ch:ltf-codes-groups}

This chapter covers three topics that connect the Boolean Fourier theory
to the thesis: linear threshold functions (the basic computational unit
of binary networks), Reed--Muller codes (the coding-theoretic
interpretation of spectral truncation), and the group-theoretic
viewpoint (which places the Walsh--Fourier basis in a broader context).

% ══════════════════════════════════════════════════════════════════════
\section{Linear Threshold Functions}\label{sec:ltf}

\begin{definition}[Linear threshold function (LTF)]\label{def:ltf}
A \emph{linear threshold function} is a Boolean function of the form
\begin{equation}\label{eq:ltf}
  f(x) = \sgn\!\left(\sum_{i=1}^n w_i x_i\right)
  = \sgn(w^\top x),
\end{equation}
where $w = (w_1, \ldots, w_n) \in \R^n$ is a weight vector and
$\sgn(z) = +1$ if $z \geq 0$ and $-1$ if $z < 0$.
\end{definition}

LTFs are also called \emph{halfspaces} or \emph{threshold gates}.
In the neural network context, a single neuron with binary inputs and a
sign activation computes exactly an LTF. A binary neural network layer
computes $d$ LTFs in parallel (one per row of the weight matrix).

\begin{example}\label{ex:ltf-examples}
\begin{itemize}
  \item \textbf{Dictator:} $f(x) = \sgn(x_1) = x_1$. Weight vector
    $w = (1, 0, \ldots, 0)$.
  \item \textbf{Majority:} $f(x) = \sgn(x_1 + \cdots + x_n)$. Weight
    vector $w = (1, 1, \ldots, 1)$.
  \item \textbf{Weighted threshold:} $f(x) = \sgn(3x_1 + x_2 - 2x_3)$.
    This function outputs $+1$ when $3x_1 + x_2 - 2x_3 > 0$.
\end{itemize}
The parity function $x_1 x_2 \cdots x_n$ is \emph{not} an LTF for
$n \geq 2$ (it cannot be computed by a single threshold gate).
\end{example}


% ──────────────────────────────────────────────────────────────────────
\subsection{Chow's Theorem}\label{sec:chow}

\begin{theorem}[Chow, 1961; {\citealp{chow1961recognition}}; see also {\citet[Chapter~5]{odonnell2014analysis}}]
\label{thm:chow}
A linear threshold function $f$ is uniquely determined among all LTFs
by its degree-0 and degree-1 Fourier coefficients:
$\fhat(\emptyset) = \E[f]$ and $\fhat(\{i\}) = \E[f \cdot x_i]$ for
$i = 1, \ldots, n$.
\end{theorem}

The $n+1$ values $(\fhat(\emptyset), \fhat(\{1\}), \ldots, \fhat(\{n\}))$
are called the \emph{Chow parameters}. Chow's theorem says: if two LTFs
have the same Chow parameters, they are the same function.

\begin{remark}\label{rem:chow-significance}
Chow's theorem implies that the ``essential information'' of an LTF lives
in degree $\leq 1$ of the Fourier spectrum. This does not mean the
higher-degree coefficients are zero---they are generally nonzero---but
that they are \emph{determined} by the low-degree ones (within the class
of LTFs). This is a strong structural result supporting the idea that
LTFs are ``low-complexity'' functions.
\end{remark}

\begin{proof}[Proof sketch]
Two distinct LTFs $f = \sgn(w^\top x)$ and $g = \sgn(v^\top x)$ define
different halfspaces in $\R^n$. If they had the same Chow parameters,
then $\E[f \cdot x_i] = \E[g \cdot x_i]$ for all $i$, meaning $f$ and
$g$ agree on average with every coordinate projection. A geometric
argument shows this forces $w$ and $v$ to define the same halfspace
(up to positive scaling), hence $f = g$.

The full proof uses the fact that an LTF is determined by the half-space
$\{z : w^\top z \geq 0\}$, and the Chow parameters encode sufficient
information about this half-space. See \citet[Chapter~5]{odonnell2014analysis}.
\end{proof}


% ──────────────────────────────────────────────────────────────────────
\subsection{Spectral Properties of LTFs}\label{sec:ltf-spectrum}

LTFs have several spectral properties that are relevant to the thesis:

\paragraph{Spectral concentration.}
LTFs are spectrally concentrated at low degrees. While the exact spectrum
depends on the weight vector $w$, the majority function (equal weights)
has 75\% of its energy at degree~1 for $n = 3$, and the concentration
improves with $n$. For general LTFs, most of the spectral energy is
within the first $O(\sqrt{n})$ degrees.

\paragraph{Noise sensitivity.}
LTFs have bounded noise sensitivity. The following result quantifies this:

\begin{theorem}[{\citet{harsha2014bounding}}]\label{thm:ltf-noise}
For any LTF $f \colon \pmone^n \to \pmone$, the total influence satisfies
$\mathbf{I}[f] = O(\sqrt{n})$.
\end{theorem}

More precisely, for $f = \sgn(w^\top x)$ with
$\norm{w}_2 = 1$: $\mathbf{I}[f] = \Theta(\norm{w}_1 / \norm{w}_2)$
(by Berry--Ess\'een arguments). For binary weights
$w \in \pmone^n$, $\norm{w}_1 = n$ and $\norm{w}_2 = \sqrt{n}$, giving
$\mathbf{I}[f] \approx \sqrt{2n/\pi}$.

\paragraph{Relevance to the thesis.}
Each row of a binary weight matrix followed by a sign activation computes
an LTF. The bounded noise sensitivity of LTFs means that:
\begin{itemize}
  \item Each neuron in a binary layer is moderately sensitive to its inputs
    (influence $O(\sqrt{d})$ out of $d$ inputs, at $d = 1024$ this is $\approx 25$).
  \item The spectral tail (high-degree coefficients) is bounded, supporting
    the spectral concentration hypothesis.
  \item The sign activation acts as a ``spectral filter,'' preserving
    low-degree structure and discarding high-frequency noise.
\end{itemize}


% ══════════════════════════════════════════════════════════════════════
\section{Reed--Muller Codes}\label{sec:reed-muller}

There is a deep connection between the Boolean Fourier basis and
Reed--Muller codes from coding theory. The thesis uses this connection
to analyze the spectral-to-binary projection.

References: \citet{reed1954class,muller1954application} for the
original constructions; \citet{macwilliams1977theory} for a comprehensive
treatment.

\subsection{Definition}\label{sec:rm-def}

\begin{definition}[Reed--Muller code]\label{def:reed-muller}
The $r$-th order \emph{Reed--Muller code} $\text{RM}(r, n)$ is the set
of all evaluation vectors of multilinear polynomials of degree $\leq r$
over $n$ Boolean variables.

Formally: identify each Boolean function $f \colon \pmone^n \to \pmone$
with its evaluation vector $v_f \in \pmone^{2^n}$. Then
\[
  \text{RM}(r, n) = \{v_f : f \text{ has Fourier degree} \leq r\}
  = \left\{\sum_{|S| \leq r} c_S\, \chi_S : c_S \in \R \text{ such that the evaluations are in } \pmone\right\}.
\]
\end{definition}

Wait---this is not quite right. The standard definition uses $\{0,1\}$
and $\text{GF}(2)$, while we are working with $\pmone$ and $\R$.
Let us give both perspectives.

\paragraph{$\text{GF}(2)$ perspective.}
Over $\text{GF}(2) = \{0, 1\}$, $\text{RM}(r, n)$ is a linear code of
length $2^n$ consisting of all evaluation vectors of polynomials
$p(b_1, \ldots, b_n)$ over $\text{GF}(2)$ of total degree $\leq r$.
Its dimension is $\sum_{i=0}^r \binom{n}{i}$ and its minimum distance
is $2^{n-r}$.

\paragraph{$\pmone$ perspective.}
Converting to $\pmone$: the evaluation vector of a degree-$\leq r$
$\text{GF}(2)$-polynomial $p$, mapped via $b \mapsto (-1)^b$, gives
a $\pmone$-vector whose Walsh--Fourier support lies in
$\{S \subseteq [n] : |S| \leq r\}$.

Conversely, any $\pmone$-vector whose Walsh--Fourier expansion has
degree $\leq r$ corresponds to a codeword in $\text{RM}(r, n)$.

\begin{theorem}[RM parameters]\label{thm:rm-params}
$\text{RM}(r, n)$ is a linear code with:
\begin{itemize}
  \item \emph{Length:} $N = 2^n$.
  \item \emph{Dimension:} $\sum_{i=0}^r \binom{n}{i}$.
  \item \emph{Minimum distance:} $2^{n-r}$.
\end{itemize}
\end{theorem}

\begin{example}\label{ex:rm-small}
\begin{itemize}
  \item $\text{RM}(0, n)$: repetition code. Only two codewords: all-$1$
    and all-$(-1)$. Minimum distance $2^n$. (Degree~0 functions are
    constants $\pm 1$.)
  \item $\text{RM}(1, n)$: the first-order RM code. Dimension $n + 1$.
    Codewords are affine functions $\pm \chi_S$ for $|S| \leq 1$ (plus
    their negations). Minimum distance $2^{n-1}$.
  \item $\text{RM}(n, n)$: the full code, containing all $2^{2^n}$
    Boolean functions. Minimum distance~1.
\end{itemize}
\end{example}


% ──────────────────────────────────────────────────────────────────────
\subsection{Connection to the Thesis}\label{sec:rm-thesis}

The thesis's degree-$k$ spectral parameterization optimizes within the
span of $\{\chi_S : |S| \leq k\}$---which is exactly the ``continuous
relaxation'' of $\text{RM}(k, n)$. The coefficients are real during
training and are projected to $\pmone$ via the sign function.

Three properties of RM codes are relevant:

\begin{enumerate}[label=(\roman*)]
  \item \textbf{Minimum distance.}
    $\text{RM}(k, n)$ has minimum distance $2^{n-k}$. For the thesis
    with $k = 3$ and $n = 10$: minimum distance $= 2^7 = 128$. This
    means any two distinct degree-$\leq 3$ Boolean functions differ on
    at least 128 of the 1024 inputs.

  \item \textbf{Dimension = parameter count.}
    The dimension $\sum_{i=0}^k \binom{n}{i}$ is exactly the number of
    spectral coefficients in the thesis's parameterization:
    $\binom{10}{\leq 3} = 1 + 10 + 45 + 120 = 176$.

  \item \textbf{Sign projection is not RM decoding.}
    The sign function applied to a continuous approximation
    $f^{\leq k}(x)$ is \emph{not} guaranteed to produce a codeword of
    $\text{RM}(k, n)$. The sign of a degree-$k$ polynomial can have
    Fourier support up to degree~$n$. The RM connection provides
    structural insight but does not directly constrain the output.
\end{enumerate}


% ══════════════════════════════════════════════════════════════════════
\section{The Group-Theoretic Perspective}\label{sec:group-theory}

The Boolean Fourier expansion is a special case of Fourier analysis on
finite abelian groups. This section briefly develops that perspective.
The main reference is \citet{terras1999fourier}; see also
\citet{diaconis1988group}.

\subsection{The Group $\Z_2^n$}\label{sec:z2n}

The Boolean hypercube $\pmone^n$ under coordinatewise multiplication is
isomorphic to the group $\Z_2^n = (\Z/2\Z)^n$ under coordinatewise
addition mod~2.

The isomorphism: $\phi \colon \Z_2^n \to \pmone^n$ by
$\phi(b_1, \ldots, b_n) = ((-1)^{b_1}, \ldots, (-1)^{b_n})$.
Then $\phi(a + b) = \phi(a) \cdot \phi(b)$ (coordinatewise).

\subsection{Characters of Finite Abelian Groups}\label{sec:characters}

\begin{definition}[Character]\label{def:character}
A \emph{character} of a finite abelian group $G$ is a group homomorphism
$\chi \colon G \to \mathbb{C}^*$ (the nonzero complex numbers under
multiplication). For $\Z_2^n$, all characters are real-valued (taking
values in $\pmone$).
\end{definition}

\begin{theorem}\label{thm:characters-z2n}
The characters of $\Z_2^n$ are exactly the functions
$\chi_S(b) = (-1)^{\sum_{i \in S} b_i}$ for $S \subseteq [n]$.
Under the $\pmone$ convention, these become $\chi_S(x) = \prod_{i \in S} x_i$:
exactly the parity functions.
\end{theorem}

\begin{proof}
A character $\chi$ of $\Z_2^n$ is determined by its values on the
standard generators $e_i = (0, \ldots, 1, \ldots, 0)$. Since
$2 e_i = 0$ in $\Z_2^n$, $\chi(e_i)^2 = \chi(0) = 1$, so
$\chi(e_i) \in \{+1, -1\}$. The $2^n$ choices for
$(\chi(e_1), \ldots, \chi(e_n)) \in \pmone^n$ give all $2^n$ characters.
\end{proof}

\begin{theorem}[Fourier analysis on finite abelian groups]\label{thm:fourier-group}
For any finite abelian group $G$ with character group $\hat{G}$, the
characters form an orthonormal basis for $L^2(G)$ (functions $G \to \mathbb{C}$
with the inner product $\ip{f}{g} = \frac{1}{|G|}\sum_{x \in G} f(x)\overline{g(x)}$).
Every function $f \colon G \to \mathbb{C}$ has a unique expansion
$f = \sum_{\chi \in \hat{G}} \fhat(\chi)\, \chi$, and Parseval holds:
$\norm{f}_2^2 = \sum_\chi |\fhat(\chi)|^2$.
\end{theorem}

For $G = \Z_2^n$, this recovers exactly the Boolean Fourier expansion of
\Cref{ch:boolean-fourier}.

\subsection{Comparison Table}\label{sec:comparison-table}

\begin{center}
\begin{tabular}{lll}
\toprule
\textbf{Classical ($\R/\Z$ or $\Z/N\Z$)} & \textbf{Boolean ($\Z_2^n$)} & \textbf{Role} \\
\midrule
$e^{2\pi i k x}$ & $\chi_S(x) = \prod_{i \in S} x_i$ & Basis function \\
Complex-valued & $\pm 1$-valued & Basis alphabet \\
DFT / FFT & WHT / FWHT & Fast algorithm \\
$O(N \log N)$ complex ops & $O(N \log N)$ integer add/sub & Complexity \\
$\sum_k |\hat{f}(k)|^2 = \norm{f}_2^2$ & $\sum_S \fhat(S)^2 = \norm{f}_2^2$ & Parseval \\
$\widehat{f \ast g}(k) = \hat{f}(k)\hat{g}(k)$ & $\widehat{f \ast g}(S) = \fhat(S)\ghat(S)$ & Convolution \\
Frequency $k$ & Degree $|S|$ & ``Frequency'' \\
Low-pass filter & Degree truncation & Smoothing \\
\bottomrule
\end{tabular}
\end{center}

The Boolean case is simpler in every respect: the basis is real (actually
$\pm 1$-valued), the transform uses only additions and subtractions, and
the group structure is particularly transparent ($\Z_2^n$ is a product of
copies of the simplest nontrivial group). This simplicity is precisely
what makes it attractive for binary neural networks.
